% \section{Beispiele}
% % Zitat ohne Seitenangabe: \cite{BoeckleKnauberPohl+04}

% % Zitat mit Seitenangabe: \cite[S.1]{BoeckleKnauberPohl+04}

% % Referenz eines Glossareintrags: \gls{computer}

% Eine normale Liste:
% \begin{itemize}
% \item Ein Punkt
% \item Ein anderer Punkt
% \end{itemize}

% Eine nummerierte Liste:
% \begin{enumerate}
% \item Erstens
% \item Zweitens
% \item \ldots
% \end{enumerate}

% Eine Tabelle:
% \begin{table}[h]
%     \centering
%     \captionabove{Eine Tabelle} 
%     \label{tab:bspTabelle}
%     \begin{tabular}{|c|c|c|} \hline
%         Spalte 1 & Spalte 2 & Spalte 3 \\ \hline
%         1 & 2 & 3 \\ \hline
%     \end{tabular}
% \end{table}

% \textbf{Fetter Text}

% \textit{Kursiver Text}

% Eine Referenz: Siehe Tabelle \ref{tab:bspTabelle}: \nameref{tab:bspTabelle}
% auf Seite \pageref{tab:bspTabelle}.

% Ein Bild:
% \begin{figure}[htb]
% \centering
% \includegraphics[width=0.2\textwidth]{Bilder/St_Uni-Logo-9-2003-eps-converted-to.pdf}
% \caption{Das Logo der SUH}
% \label{fig:suh-logo}
% \end{figure}

% Eine abgesetzte Formel:
% \begin{equation}
% \sum_{n=0}^{3}n=6
% \end{equation}

% Eine Formel im Flie�text: $2^2=4$ 

% Ein Listing aus einer Datei:
% \lstset{language=Java, basicstyle=\footnotesize, showstringspaces=false,tabsize=2}
% \lstinputlisting[label=lst:HelloWorld,caption={HelloWorld}]{Listings/HelloWorld.java}

% Ein 'on-the-fly' erstelltes Listing:
% \begin{lstlisting}[label=lst:log,caption={Beispiel eines Log-Eintrags}] 
% (Sun Sep 13 23:02:20 2009): ODBC Driver <system32>\wbemdr32.dll not present 
% (Sun Sep 13 23:02:20 2009): Successfully verified WBEM OBDC adapter (incompatible version removed if it was detected).
% (Sun Sep 13 23:02:20 2009): Wbemupgd.dll Registration completed.
% \end{lstlisting}


